\phantomsection\addcontentsline{toc}{chapter}{课程说明与日程}
\markboth{课程说明与日程}{}
\noindent{\color{structurecolor}\bfseries\fontsize{14.4}{20}\selectfont 课程说明与日程}\par
本节依两门课程的官方 syllabus、calendar、assignments 与 exams 网页译编, 保留历史课程的安排, 不将其解释为本书读者的考核规定. 网页快照及确切下载地址随源码保存. 第1--14章构成已有主线与拓展正文; 后续章节按原编号收入160道公开题卷大题及解答, 包括361个印刷叶项, 其中1项是阅读与流线观察任务, 其余360项为数学解答任务. 另有原14章的38道自设练习. 这些数按题目出现次数统计, 不把重复题当成新的数学结论.

\noindent\textbf{18.04: Complex Variables with Applications, Spring 2018.}
授课教师 Jeremy Orloff. 课程讲义的设计源于 Andre Nachbin, Jörn Dunkel 参与校订; 本次收入的习题课题面、答案页内署名为 Vishesh Jain. 每周3次讲课, 每次1小时; 每周1次习题课, 1小时. 官方先修是18.02多元微积分, 加18.03或18.032微分方程. 无指定必需教材, 主要读课程讲义; Michael Taylor 的 \textit{Introduction to Complex Analysis}, Beck、Marchesi、Pixton、Sabalka 的 \textit{A First Course in Complex Analysis} 与 Brown、Churchill 第9版教材为补充阅读, 本档案不收入这些第三方书.

课程主线为解析函数、Cauchy 定理与积分公式、Taylor 和 Laurent 级数; 应用包括调和函数、二维流体势、实积分、保角映射、辐角原理、Laplace 和 Fourier 变换, 最后涉及解析延拓及 Gamma 函数. 9次作业, 两次1小时课内考试, 一次3小时期末考试. 作业占25\%, 两次课内考试合占40\%, 期末占35\%; 期末必须及格才能通过课程. 允许讨论合作, 但提交解答须以自己的语言书写并列出合作人及外部来源. 公开 exams 网页提供三份\emph{模拟题}, 不应冒称它们是2018实际考试卷. 公开 recitations 网页无第10次题卷或答案, 收入的是1--9、11--13这12次.

\begin{longtable}{@{}p{.09\textwidth}p{.64\textwidth}p{.21\textwidth}@{}}
\toprule 课次&主题（保留原 Topic 编号）&关键安排\\\midrule\endhead
1&Topic 1: 复数代数与复平面&\\
2&Topic 1续&\\
3&Topic 2: 解析函数&作业1截止\\
4&Topic 2续&\\
5&Topic 2续&\\
6&Topic 2续&作业2截止\\
7&18.02多元微积分复习&\\
8&Topic 3: Cauchy 定理&\\
9&Topic 3续&作业3截止\\
10&Topic 4: Cauchy 积分公式&\\
11&Topic 4续&\\
12&Topic 4续&作业4截止\\
13&Topic 5: 调和函数&\\
14&第一次考试复习&\\
15&第一次考试, Topics 1--4&考试1\\
16&Topic 5续; Topic 6: 二维流体力学与复势&\\
17&Topic 6续&\\
18&Topic 7: Taylor 与 Laurent 级数&作业5截止\\
19&Topic 7续&\\
20&Topic 7续&\\
21&Topic 8: 留数定理&\\
22&Topic 8续&\\
23&Topic 8续&\\
24&Topic 9: 用留数计算定积分&作业6截止\\
25&Topic 9续&\\
26&Topic 9续&\\
27&Topic 10: 保角变换&作业7截止\\
28&第二次考试复习&\\
29&第二次考试, Topics 5--9&考试2\\
30&Topic 10续&\\
31&Topic 10续&\\
32&Topic 11: 辐角原理&作业8截止\\
33&Topic 11续&\\
34&Topic 11续&\\
35&Topic 11续&\\
36&Topic 12: Laplace 变换&作业9截止\\
37&Topic 12续&\\
38&Topic 12续&\\
39&Topic 13: 解析延拓与 Gamma 函数&\\
40&期末考试&期末\\\bottomrule
\end{longtable}

\noindent\textbf{18.112: Functions of a Complex Variable, Fall 2008.}
教师 Sigurdur Helgason; 补充讲义由 Zuoqin Wang 在其指导下编写. 每周2次讲课, 每次1.5小时. 先修18.100B（Analysis I）或同等课程. 指定 Lars V. Ahlfors 的 \textit{Complex Analysis}, 第3版, McGraw-Hill, 1979, ISBN 9780070006577. 教师计划讲授至第232页的大部分内容, 另加两讲完整证明素数定理; 公开补充讲义并不是指定书的完整替代品. 本书选择其中与复分析主线相关的拓展; 素数定理和 zeta 函数函数方程的完整论证仍属未覆盖内容, 见来源表.

每隔一周星期四布置作业, 下一个星期四截止, 从 Ahlfors 书中选题. 网页列有六批共24项页码/题号及 MIT 公开解答, 没有原书完整题面, 也没有逐批的具体公历日期. 因此本书准确登记引用和解答文件, 不从解答反推并伪称翻译了原书题目. 两次课内考试分别在 L4、L16 之后, 另有期末; exams 网页说明均开卷. 作业10\%, 两次课内考试各20\%（合40\%）, 期末50\%, 以累计分数评定. 可与他人讨论, 但需给自己的整洁解答. 此归档页面链接的三份卷面日期实际为2006年10月20日、11月22日、12月20日, 外加2008课程封面; 下文明确标为2006历史题卷, 不改题卷年份.

\begin{longtable}{@{}p{.11\textwidth}p{.83\textwidth}@{}}
\toprule 课次&主题\\\midrule\endhead
L1&复数代数、复平面几何、球面表示\\
L2&复指数与复对数; 指数、三角函数与对数分支\\
L3&解析函数和有理函数; Cauchy--Riemann 方程\\
L4&复幂级数与一致收敛\\
E1&第一次课内考试\\
L5&指数与三角函数\\
L6&保角映射、线性变换; 解析函数的局部几何、保角性与尺度变换\\
L7&线性变换续; 交比、对称性、圆的作用\\
L8&线积分; 路径无关与原函数存在的等价性\\
L9&Cauchy--Goursat 定理\\
L10&特殊 Cauchy 公式及应用; 可去奇点、有余项的复 Taylor 定理\\
L11&孤立奇点\\
L12&局部映射、Schwarz 引理与非欧解释; 拓扑性质、最大模定理\\
L13&一般 Cauchy 定理\\
L14&留数定理及应用; 留数计算、辐角原理、Rouché 定理\\
L15&围道积分及应用; 定积分及对数分支的处理\\
L16&调和函数及与全纯函数的关系; Poisson 公式、Schwarz 定理\\
E2&第二次课内考试\\
L17&Mittag-Leffler 定理; Laurent 级数与部分分式展开\\
L18&无穷乘积、Weierstrass 典范乘积、Gamma 函数\\
L19&正规族; 全纯函数的等度有界性、Arzelà 定理\\
L20&Riemann 映射定理\\
L21--L22&素数定理: 历史、完整证明及证明细节（官方日程合列两讲）\\
L23&zeta 函数延拓到复平面及函数方程\\
E3&期末考试\\\bottomrule
\end{longtable}
\noindent 依据: \href{https://ocw.mit.edu/courses/18-04-complex-variables-with-applications-spring-2018/pages/syllabus/}{18.04课程说明}, \href{https://ocw.mit.edu/courses/18-04-complex-variables-with-applications-spring-2018/pages/calendar/}{18.04日程}, \href{https://ocw.mit.edu/courses/18-112-functions-of-a-complex-variable-fall-2008/pages/syllabus/}{18.112课程说明}, \href{https://ocw.mit.edu/courses/18-112-functions-of-a-complex-variable-fall-2008/pages/calendar/}{18.112日程}. 2026-10-08核查.
